\section{Noetherian rings}
\lecture\label{lec:NoetherianRings}

\begin{definition}
\index{Ring!noetherian}
	A ring $R$ is said to be \emph{noetherian} if every (increasing)
	sequence \[
 I_1\subseteq I_2\subseteq\cdots
 \]
 of ideals of $R$
	stabilizes, that is $I_n=I_m$ for some $m\in\Z_{>0}$ and all $n\geq m$. 
\end{definition}

Finite rings
are noetherian. The ring $\Z$ of integers is noetherian. Why?

\begin{exercise}
Let $R=\{f\colon [0,1]\to\R\}$ with 
\[
(f+g)(x)=f(x)+g(x),
\quad
(fg)(x)=f(x)g(x),
\quad
f,g\in R,\,x\in [0,1].
\]
For $n\in\Z_{>0}$ let
$I_n=\{f\in R:f|_{[0,1/n]}=0\}$. Prove that each $I_n$ is an ideal of $R$ and 
the sequence 
$I_1\subsetneq I_2\subsetneq\cdots$ 
does not stabilize. Conclude that $R$ is not noetherian. 
\end{exercise}

\begin{definition}
\index{Ideal!finitely generated}
	Let $R$ be a ring. An ideal $I$ of $R$ 
	is said to be \emph{finitely generated} if $I=(X)$ for some
	finite subset $X$ of $R$. 
\end{definition}

If $R$ is a ring, $\{0\}$ and $R$ are finitely generated. Why?

If $X=\{x_1,\dots,x_n\}$ is a subset of a ring $R$, 
we write $(x_1,\dots,x_n)$ to denote 
the ideal of $R$ generated by the set $\{x_1,\dots,x_n\}$. 

\begin{proposition}
Let $R$ be a ring. Then $R$ is noetherian if and only 
if every ideal of $R$ is finitely generated. 	
\end{proposition}

\begin{proof}
  Assume first that $R$ is noetherian. Let $I$ be an ideal of $R$ that is not
  finitely generated.  Thus $I\ne\{0\}$. Let $x_1\in I\setminus\{0\}$ and let
  $I_1=(x_1)$. Since $I$ is not finitely generated, $I\ne I_1$ and hence
  $\{0\}\subsetneq I_1\subsetneq I$. Let $x_2\in I\setminus I_1$. Then
  $I_2=(x_1,x_2)$ is a finitely generated ideal such that $I_1\subsetneq
  I_2\subsetneq I$. We continue with this procedure. Once we have the ideals
  $I_1,\dots,I_{k-1}$, let $x_k\in I\setminus I_{k-1}$ (such an element exists
  because $I_{k-1}$ is finitely generated and $I$ is not) and
  $I_k=(I_{k-1},x_k)$. The sequence $\{0\}\subsetneq I_1\subsetneq
  I_2\subsetneq\cdots$ does not stabilize, a contradiction to $R$ being 
  noetherian. 
	
	Assume now that every ideal of $R$ is finitely generated and 
	let $I_1\subseteq I_2\subseteq\cdots$ be a sequence of ideals of $R$. Then
	$I=\bigcup_{i\geq1}I_i$ is an ideal of $R$, so it is finitely generated, say
	$I=(x_1,\dots,x_n)$. We know that $x_j\in I_{i_j}$ for all $j$. Let 
	$N=\max\{i_1,\dots,i_n\}$. Then 
	$x_j\in I_N$ for all $j\in\{1,\dots,n\}$ 
	and hence $I\subseteq I_N$. This implies that 
	$I_m=I_N$ for all $m\geq N$. 
\end{proof}

\begin{exercise}
	Let $R=\C[X_1,X_2,\dots]$ be the ring of polynomials in 
  countably infinitely many commuting
  variables. 
	Every element of $R$ is a finite sum
	of the form 
	\[
	\sum a_{i_1i_2\cdots i_k}X_{i_1}^{n_1}X_{i_2}^{n_2}\cdots X_{i_k}^{n_k}
	\]
	for some complex numbers 
	$a_{i_1i_2\cdots i_k}$. In particular, 
  for every $f\in R$ there exists $n\geq1$
  such that
  \[
  f\in\C[X_1,\dots,X_n].
  \]
	Prove that the ideal $I=(X_1,X_2,\dots)$ of polynomials 
	with zero constant term is not finitely generated. 
\end{exercise}

The correspondence theorem and the previous proposition 
allow us to prove easily the following result. 

\begin{proposition}
	Let $R$ be a ring and
  $I$ an ideal of $R$. If $R$ is noetherian, then $R/I$ is noetherian.
\end{proposition}

\begin{proof}
	Let $\pi\colon R\to R/I$ be the canonical surjection and let $J$ be an ideal of $R/I$. 
	Then $\pi^{-1}(J)$ is an ideal of $R$ containing $I$. Since 
	$R$ is noetherian, $\pi^{-1}(J)$ is finitely generated, say 
	$\pi^{-1}(J)=(x_1,\dots,x_k)$ for $x_1,\dots,x_k\in R$. Thus 
	\[
	J=\pi(\pi^{-1}(J))=(\pi(x_1),\dots,\pi(x_k)),
	\]
	because $\pi$ is surjective 
	and hence $J$ is finitely generated. 
\end{proof}

Since $\Z$ is noetherian, $\Z/n$ is noetherian for all $n\geq2$. 

\begin{exercise}
\label{xca:R[X]_noetherian}
	Prove that $\R[X]$ is noetherian. 	
\end{exercise}

Hint: Use the fact that every ideal of $\R[X]$ is principal
(Exercise~\ref{xca:R[X]_principal}).  

\section{Hilbert's basis theorem}

\begin{theorem}[Hilbert]
\index{Hilbert's basis theorem}
	Let $R$ be a commutative ring. If $R$ is a noetherian ring, then $R[X]$ is noetherian.	
\end{theorem}

\begin{proof}
	We must show that every ideal of $R[X]$ is finitely generated. Assume that
	there is an ideal $I$ of $R[X]$ that is not finitely generated. In particular, $I\ne\{0\}$. 
	Let $f_1(X)\in I\setminus\{0\}$ be of minimal degree $n_1$. 
	Since $I$ is not finitely generated, it follows that 
	$I\ne (f_1(X))$. Let $f_2(X)\in I\setminus (f_1(X))$ be
	of minimal degree $n_2$. In particular, the minimality of the degree of $f_1(X)$ implies that 
	$n_2\geq n_1$. 
	We continue with this procedure. For $i>1$ let 
	$f_i(X)\in I$ be a polynomial of minimal degree $n_i$   
	such that $f_i(X)\not\in(f_1(X),\dots,f_{i-1}(X))$ (note
	that such an $f_i(X)$ exists because $I$ is not finitely generated). 
	Moreover, $n_i\geq n_{i-1}$. This happens because 
	if $n_i<n_{i-1}$, then
	$f_i(X)\not\in (f_1(X),\dots,f_{i-2}(X))$, which contradicts
	the minimality of $n_{i-1}=\deg f_{i-1}(X)$. 

	For each $i\geq1$ 
	let $a_i$ be the leading coefficient of $f_i(X)$, that is
	\[
	f_i(X)=a_iX^{n_i}+\cdots,
	\]
	where the dots denote a polynomial of 
	degree $<n_i$. Note that 
	$a_i\ne 0$ for all $i\geq 1$. 
	
	Let $J=(a_1,a_2,\dots)$. Since $R$ is noetherian, the sequence
	\[
	(a_1)\subseteq (a_1,a_2)\subseteq\cdots\subseteq(a_1,a_2,\dots,a_k)\subseteq\cdots
	\]
	stabilizes, so we may assume that 
	$J=(a_1,\dots,a_m)$ for some $m\in\Z_{>0}$. 
	In particular, there exist $u_1,\dots,u_m\in R$ such that 
	\[
	a_{m+1}=\sum_{i=1}^m u_ia_i.
	\]
	Let 
	\[
	g(X)=\sum_{i=1}^mu_if_i(X)X^{n_{m+1}-n_i}\in (f_1(X),\dots,f_m(X))\subseteq I.
	\]
	The leading coefficient of $g(X)$ is $\sum_{i=1}^mu_ia_i=a_{m+1}$ and, moreover, 
	the degree of $g(X)$ is $n_{m+1}$. Thus $\deg(g(X)-f_{m+1}(X))<n_{m+1}$. 
	
	Since $f_{m+1}(X)\not\in (f_1(X),\dots,f_m(X))$, 
	\[
	g(X)-f_{m+1}(X)\not\in (f_1(X),\dots,f_m(X)),
	\]
	a contradiction to the minimality of the degree of $f_{m+1}$.  
\end{proof}

Since $R[X_1,\dots,X_n]=(R[X_1,\dots,X_{n-1}])[X_n]$, by induction 
one proves that if $R$ is a commutative noetherian ring, 
then $R[X_1,\dots,X_n]$ is noetherian. 
 
\begin{example}
    Let $N\ne1$ be a square-free integer.  
	Since $\Z$ is noetherian, so is $\Z[X]$ by Hilbert's theorem. Now 
	$\Z[\sqrt{N}]$ is noetherian, as \[
    \Z[\sqrt{N}]\simeq\Z[X]/(X^2-N)
    \]
    and quotients
	of noetherian rings are noetherian.  	
\end{example}

\begin{example}
	The ring $\Z[X,X^{-1}]$ is noetherian, as 
    \[
    \Z[X,X^{-1}]\simeq\Z[X,Y]/(XY-1).
    \]
\end{example}
 

\begin{exercise}
\label{xca:Hilbert_powerseries}
\index{Ring!of formal power series}
   Let $R$ be a commutative ring and $R[\![X]\!]$ be the ring of formal power series with the usual operations.
   Prove that $R[\![X]\!]$ is noetherian if $R$ is noetherian.
\end{exercise}

\begin{exercise}
	Let $f\colon R\to R$ be a surjective ring homomorphism. Prove that $f$ is an isomorphism
	if $R$ is noetherian. 	
\end{exercise}


